% Source: https://artofproblemsolving.com/community/c3h1349710_are_all_powers_of_1_one
% Stable identity: aops-topic-1349710-q01
(a) The base case and the step $n\to n+1$ prove $1^n=1$ for every \emph{positive integer} $n$, and nothing more: induction reaches exactly the numbers $1,2,3,\dots$ that can be obtained from $1$ by repeatedly adding $1$. The argument is correct for these.

(b) The number $\pi$ is not a positive integer, so it is never reached by the induction; moreover, for a non-integer exponent $1^{\pi}$ is not defined by repeated multiplication at all, so the step $1^{n+1}=1^n\cdot1$ says nothing about it. A correct argument has to use the definition of a real power:
\begin{itemize}
\item For a rational exponent $\dfrac{p}{q}$ with $q>0$, $1^{\dfrac{p}{q}}=\left(\sqrt[q]{1}\right)^p=1^p=1$, since the positive $q$-th root of $1$ is $1$ (and $1^{-p}=\dfrac{1}{1^p}=1$).
\item For a real exponent $r$, $a^r$ is defined (for $a>0$) as the limit of $a^{r_n}$ for rational $r_n\to r$, equivalently $a^r=e^{r\ln a}$. With $a=1$ every $1^{r_n}$ equals $1$, so the limit is $1$; or directly $1^{\pi}=e^{\pi\ln1}=e^{0}=1$.
\end{itemize}
So $1^{\pi}=\boxed{1}$, but this needs the definition of irrational powers, not induction.
